\documentclass[11pt, a4paper]{article}

% ── Packages ──────────────────────────────────────────────────────────────────
\usepackage[margin=2.5cm]{geometry}
\usepackage{amsmath, amssymb}
\usepackage{booktabs}
\usepackage{array}
\usepackage{tabularx}
\usepackage{hyperref}
\usepackage{natbib}
\usepackage{setspace}
\usepackage{titlesec}
\usepackage{parskip}
\usepackage{xcolor}
\usepackage{fancyhdr}

% ── Colors ────────────────────────────────────────────────────────────────────
\definecolor{ieblue}{RGB}{13, 27, 94}
\definecolor{iegreen}{RGB}{126, 200, 58}

% ── Section formatting ────────────────────────────────────────────────────────
\titleformat{\section}{\large\bfseries\color{ieblue}}{}{0em}{}[\vspace{-4pt}\rule{\linewidth}{1pt}]
\titleformat{\subsection}{\normalsize\bfseries\color{ieblue}}{}{0em}{}

% ── Header / Footer ───────────────────────────────────────────────────────────
\pagestyle{fancy}
\fancyhf{}
\rhead{\small\color{gray} Thesis Proposal · Yaxin Wu · IE University}
\rfoot{\small\thepage}
\renewcommand{\headrulewidth}{0pt}

% ── Hyperref ──────────────────────────────────────────────────────────────────
\hypersetup{
  colorlinks=true,
  linkcolor=ieblue,
  citecolor=ieblue,
  urlcolor=ieblue
}

% ─────────────────────────────────────────────────────────────────────────────
\begin{document}

% ── Title block ───────────────────────────────────────────────────────────────
\begin{center}
  {\LARGE\bfseries\color{ieblue}
    Carbon-Intelligent Workload Scheduling\\[4pt]
    for AI Data Centers}\\[10pt]
  {\large A Deterministic Mathematical Optimization Approach}\\[18pt]
  \begin{tabular}{rl}
    \textbf{Student}     & Yaxin Wu \\
    \textbf{Supervisor}  & Prof.\ Bissan Ghaddar \\
    \textbf{Program}     & Master of Business Analytics \& Big Data, IE University \\
    \textbf{Date}        & May 2026 \\
  \end{tabular}
\end{center}

\vspace{6pt}
\noindent\rule{\linewidth}{1.5pt}
\vspace{4pt}

% ── 1. Research Positioning ───────────────────────────────────────────────────
\section{Research Positioning}

\textbf{Supervisor's framework}: How to build a system that translates grid carbon intensity signals into real-time data center scheduling decisions, using Virtual Capacity Curves (VCC) derived from ElectricityMaps data and Google Cluster Traces \citep{colangelo2025}.

\textbf{This thesis}: Within that framework, how can a \emph{deterministic mathematical programming model}, rather than heuristic rules, minimize AI data center carbon emissions through joint temporal and spatial workload scheduling across five grid regions?

The supervisor defines the overall pipeline (VCC generation, workload classification, simulation architecture). This thesis formalizes the Carbon-Aware Optimization component into a provably optimal LP, deepening and extending the existing framework.

% ── 2. Core Research Question ─────────────────────────────────────────────────
\section{Core Research Question}

\begin{quote}
  \itshape
  How can a deterministic mathematical programming model, using carbon intensity (CI) and renewable energy fraction (RF) as scheduling signals, minimize AI data center carbon emissions through joint temporal and spatial workload scheduling across five grid regions?
\end{quote}

% ── 3. Data ───────────────────────────────────────────────────────────────────
\section{Data}

\subsection{Sources and Coverage}

\begin{table}[h]
\centering\small
\renewcommand{\arraystretch}{1.3}
\begin{tabularx}{\linewidth}{p{4.8cm}Xp{1.9cm}p{1.7cm}}
\toprule
\textbf{Signal} & \textbf{Source} & \textbf{Granularity} & \textbf{Period} \\
\midrule
Carbon Intensity CI (gCO\textsubscript{2}eq/kWh) & ElectricityMaps API & Hourly & 2024--2025 \\
Renewable Fraction RF (\%) & ElectricityMaps API & Hourly & 2024--2025 \\
Electricity Price (constraint) & PJM / NYISO / ENTSO-E / EMC & Hourly & 2024--2025 \\
Workload traces & Google Cluster Traces & Per-job & -- \\
\bottomrule
\end{tabularx}
\caption{Data sources and coverage. CI and RF are retrieved via the ElectricityMaps
\texttt{carbonIntensity} and \texttt{powerBreakdown} API fields respectively.}
\end{table}

\subsection{Target Regions}

\begin{table}[h]
\centering\small
\renewcommand{\arraystretch}{1.3}
\begin{tabularx}{\linewidth}{p{2.9cm}p{3.8cm}p{3.0cm}X}
\toprule
\textbf{Zone} & \textbf{Grid mix} & \textbf{CI level} & \textbf{Scheduling opportunity} \\
\midrule
US-MIDA-PJM & Coal / gas / nuclear & High, high variance & Highest \\
US-NY-NYIS  & Hydro + nuclear      & Medium, intra-day variability & High \\
FI          & Hydro + nuclear + wind & Low, strong seasonal variation & Medium \\
BE          & Nuclear, EU interconnect & Medium-low, neighbour influence & Medium \\
SG          & Almost entirely gas  & High but stable, low variance & Lowest \\
\bottomrule
\end{tabularx}
\caption{Five target regions. SG is an important control case: its low CI variability
means temporal shifting yields limited savings, validating the argument that not all
regions benefit equally from carbon-aware scheduling.}
\end{table}

% ── 4. Indicators ─────────────────────────────────────────────────────────────
\section{Core Indicators}

\paragraph{Carbon Intensity (CI).}
The primary optimization signal, measured in gCO\textsubscript{2}eq/kWh and fully comparable across regions. Retrieved directly from the ElectricityMaps \texttt{carbonIntensity} field \citep{radovanovic2023}.

\paragraph{Renewable Fraction (RF).}
RF measures the supply-side structure of the grid: the share of generation from variable renewables (wind, solar, hydro), computed from the ElectricityMaps \texttt{powerBreakdown} field \citep{radovanovic2023}. Two hours with identical CI can have very different RF: one driven by wind surplus, one by nuclear baseload. RF discriminates between these cases, providing a complementary signal to CI. The specific form in which RF enters the objective is derived from OLS regression of CI on RF across all five regions (EDA, Section~10); the multiplicative discount $\mathrm{CI}_r(t)\cdot(1-\alpha\cdot\mathrm{RF}_r(t))$ is supported by the finding that the normalized slope $\beta_1/\overline{\mathrm{CI}}$ is more consistent across regions than the raw slope \citep{liu2012}.

\paragraph{Electricity Price.}
Treated as a per-region budget constraint rather than an objective term, preserving model generalizability across electricity markets with incomparable absolute prices (PJM LMP, Nord Pool, USEP). Deep integration with market mechanisms is left as future work.

% ── 5. Methodology ────────────────────────────────────────────────────────────
\section{Methodology}

The core methodology is deterministic optimization. Historical CI and RF data are used directly as model inputs, and evaluation is performed via backtesting over a two-year historical window.

{\small\[
  \underbrace{\text{Historical CI/RF (ElectricityMaps)}}_{\text{Input}}
  \;\longrightarrow\;
  \underbrace{\text{Deterministic Optimization (Gurobi LP)}}_{\text{Core contribution}}
  \;\longrightarrow\;
  \underbrace{\text{Backtesting vs.\ baselines}}_{\text{Evaluation}}
\]}

\subsection{Unified Joint Optimization Model}

Rather than separating temporal and spatial shifting into sequential phases, the model optimizes both dimensions simultaneously. The decision variable \(x_{r,t} \geq 0\) denotes the compute load (kW) assigned to region \(r \in \mathcal{R}\) at hour \(t \in \mathcal{T}\), and covers both the \emph{when} and the \emph{where} of scheduling in a single formulation.

\subsubsection*{Objective Function}

\begin{equation}
  \min \sum_{r \in \mathcal{R}} \sum_{t \in \mathcal{T}}
    x_{r,t} \cdot \text{CI}_r(t) \cdot \bigl(1 - \alpha \cdot \text{RF}_r(t)\bigr)
  \label{eq:objective}
\end{equation}

Both CI and RF are deterministic input parameters: hourly values retrieved from the ElectricityMaps API before the optimization runs. Neither is a decision variable. They represent two distinct but complementary dimensions of grid cleanliness. CI is the direct carbon outcome: it measures actual emissions per kWh consumed and is the primary optimization signal~\citep{radovanovic2023}. RF measures the supply-side structure: the share of generation coming from variable renewables (wind, solar, hydro). A low-CI hour driven by wind and a low-CI hour driven by nuclear have the same CI value but different scheduling implications; RF discriminates between them.

The motivation for including RF as an independent signal draws on \citet{radovanovic2023}, who argue for hourly carbon-free energy tracking that distinguishes renewable and nuclear contributions, and on \citet{liu2012}, who establish precedent for incorporating renewable fraction into LP scheduling objectives. The specific multiplicative form $\mathrm{CI}_r(t)\cdot(1-\alpha\cdot\mathrm{RF}_r(t))$ is an original construction of this thesis, grounded in the EDA regression analysis (Section~10). OLS regression of CI on RF across all five regions yields a statistically significant negative slope ($p \ll 0.01$ in every region), and the normalized slope $\beta_1/\overline{\mathrm{CI}}$ is more consistent across regions than the raw slope $\beta_1$. This indicates that the RF effect scales proportionally with the ambient CI level rather than acting as a fixed absolute reduction, which is precisely what the multiplicative discount encodes. A single shared $\alpha \in [0,1]$ is used if the implied per-region values $\alpha_r \approx -\beta_1/\beta_0$ cluster closely; otherwise region-specific values are applied. For any region where the CI--RF correlation is weak ($|r| < 0.2$), $\alpha_r$ is set to zero.

\subsubsection*{Constraints}

\begin{align}
  &\sum_{r \in \mathcal{R}} \sum_{t \in \mathcal{T}} x_{r,t} = D
    \tag{C1}\label{eq:c1} \\[4pt]
  &x_{r,t} \leq C_{\max,r} \quad \forall\, r, t
    \tag{C2}\label{eq:c2} \\[4pt]
  &\sum_{t \in \mathcal{T}} x_{r,t} \cdot P_r(t) \leq B_r \quad \forall\, r
    \tag{C3}\label{eq:c3} \\[4pt]
  &x_{r,t} \geq 0 \quad \forall\, r, t
    \tag{C4}\label{eq:c4}
\end{align}

C1 enforces exact demand completion. C2 bounds regional server capacity. C3 caps per-region electricity expenditure: since prices across PJM, NYISO, Nord Pool, and USEP are denominated in different currencies, a budget cap preserves model generalizability while still bounding operational cost; electricity price data collection is in progress and C3 will be activated once sources are confirmed. C4 is non-negativity.

The supervisor's VCC mechanism \citep{colangelo2025}, which indexes flexible load capacity against a CI ceiling, is a planned extension once Google Cluster Trace data is available. It will be incorporated as an additional constraint on the flexible-load component of $x_{r,t}$.

\subsubsection*{Model Validation: Component-Level References}

\begin{table}[h]
\centering
\begin{tabularx}{\linewidth}{p{3.6cm}Xp{4cm}}
\toprule
\textbf{Component} & \textbf{Description} & \textbf{Reference} \\
\midrule
Objective form \eqref{eq:objective}
  & CI-weighted load as carbon cost signal
  & \citet{radovanovic2023} \\
RF weighting \((1-\alpha\cdot\text{RF})\)
  & Multiplicative discount derived from EDA regression; motivation from \citet{radovanovic2023} and \citet{liu2012}; functional form original to this thesis
  & This thesis \\
C1--C2
  & Demand completion and capacity bounds
  & \citet{liu2012} \\
C3 (budget)
  & Per-region electricity cost cap (pending price data)
  & \citet{hao2024joint} \\
Multi-region joint LP structure
  & Cross-region temporal-spatial allocation
  & \citet{riepin2024spatiotemporal}; \citet{attenni2024shifting} \\
\bottomrule
\end{tabularx}
\caption{Component-level references validating the mathematical formulation.}
\end{table}

\subsubsection*{Solver and Heuristic Fallback}
The model is solved with \textbf{Gurobi} as a Linear Program (LP) over continuous variables. If workload discretization is required, the formulation upgrades to a Mixed Integer Program (MIP); given the large-scale horizon (\(|\mathcal{R}|=5\), \(|\mathcal{T}| = 17{,}544\) hours over two years), LP relaxation is the default.

If the LP solve time proves prohibitive for rolling-window backtesting, a \textbf{greedy heuristic} is implemented as a fallback: at each decision window, sort hours by effective carbon cost \(\text{CI}_r(t)\cdot(1-\alpha\cdot\text{RF}_r(t))\) and greedily assign load to the cheapest region-hour pairs subject to capacity constraints. The heuristic provides a near-optimal bound and serves as an additional baseline for benchmarking solution quality.

\subsection{Evaluation: Backtesting}

The optimal schedule is evaluated by replaying two years of historical data in a rolling-window fashion, simulating real-time decisions without look-ahead beyond the current optimization window.

\begin{table}[h]
\centering
\begin{tabularx}{\linewidth}{p{3cm}X}
\toprule
\textbf{Baseline} & \textbf{Description} \\
\midrule
Carbon-Agnostic     & Uniform scheduling; no carbon awareness \\
Simple Shifting     & Rule-based: defer all flexible load to the lowest-CI hour \\
Oracle              & Perfect 24h CI foresight (theoretical upper bound) \\
\textbf{This model} & Deterministic LP with historical CI/RF (constraints C1--C4) \\
\bottomrule
\end{tabularx}
\caption{Evaluation baselines.}
\end{table}

Carbon savings are reported per region against all three baselines. A post-hoc counterfactual decomposition then isolates the marginal contribution of temporal versus spatial flexibility. Finally, CV is used to explain why savings vary across regions: high-CV grids such as PJM benefit substantially from shifting, while low-CV grids such as Singapore do not.

% ── 6. Thesis Structure ───────────────────────────────────────────────────────
\section{Planned Thesis Structure (\textasciitilde 30 Pages)}

\begin{enumerate}
  \item \textbf{Introduction} (\textasciitilde 4 pages)\\
  Establishes the motivation: AI data centers now draw hundreds of megawatts continuously, yet most scheduling systems ignore the hourly variation in grid carbon intensity that makes shifting valuable. The section sets out the research question and the three contributions, and explains why a deterministic LP fills a gap that existing heuristic approaches leave open.

  \item \textbf{Literature Review} (\textasciitilde 6 pages)\\
  Review of the carbon-aware workload scheduling literature, covering temporal shifting, spatial load balancing, joint spatio-temporal optimization, and AI-specific workload management. The review closes by positioning this thesis relative to prior work: the key gap is the absence of a deterministic LP that jointly optimizes both scheduling dimensions with an RF-weighted objective, evaluated on real multi-region data.

  \item \textbf{Problem Description} (\textasciitilde 5 pages)\\
  States the problem formally and describes the data: two years of hourly CI and RF from ElectricityMaps across the five regions. Defines CI, RF, and CV, then presents the unified joint LP, covering the objective function, constraints C1--C5, and notation.

  \item \textbf{Computation Results \& Solution Methodology} (\textasciitilde 10 pages)\\
  Solver: Gurobi LP, with a greedy carbon-aware heuristic as fallback if rolling-window solve times prove prohibitive. Results are organized as follows: single-region temporal optimization across all five zones, followed by the full joint multi-region LP. Backtesting compares the model against the three baselines. The analysis concludes with a post-hoc decomposition separating the marginal contributions of temporal and spatial flexibility, and a CV-based regression explaining cross-region variation.

  \item \textbf{Conclusion} (\textasciitilde 3 pages)\\
  Summarizes the main findings and their practical implications for data center operators. Discusses limitations of the deterministic approach and outlines directions for future work, including stochastic extensions and real-time deployment.
\end{enumerate}

% ── 7. Expected Contributions ─────────────────────────────────────────────────
\section{Expected Contributions}

The primary technical contribution is a unified joint LP that optimizes temporal and spatial workload scheduling in a single pass, without pre-imposing a sequential decomposition. The RF weighting term extends prior CI-only objectives~\citep{liu2012, radovanovic2023} by incorporating supply-side renewable structure: a low-CI window driven by wind receives a stronger incentive than one driven by nuclear. The specific multiplicative discount form is derived from EDA regression analysis rather than assumed a priori. The supervisor's VCC mechanism \citep{colangelo2025} is incorporated as a planned constraint extension once the required workload trace data is available.

On the empirical side, the model is evaluated across five structurally distinct grid regions spanning the US, Northern Europe, and Southeast Asia. CV is used as a diagnostic to explain why savings differ across regions, and the post-hoc decomposition quantifies how much of the total gain comes from shifting \emph{when} versus \emph{where} workloads run.

% ── References ────────────────────────────────────────────────────────────────
\bibliographystyle{plainnat}
\bibliography{references}

\end{document}
